\section{Introduction}
\label{sec:introduction}

Graph interventions are often justified by the statistics they preserve.
A rewiring may keep every degree, and a motif deletion may be paired with a
count of the circuit family that remains. These statistics describe the edit.
They do not, by themselves, describe what reaches the readout of a model.

A MaleCNS intervention makes the distinction concrete. After compensating
swaps, two reconstructed graphs have the same in-degree and out-degree
sequences, the same module-pair counts, and zero common-relay triples. One
still carries both designated sources to the sinks in two propagation steps;
the other carries only one. Their archived accuracies are 0.9745 and 0.5055.
A relay count alone cannot explain this difference: separate relays can
replace a deleted common relay.

We audit the information that survives an intervention for a fixed executed
operator. Features undergo two rounds of incoming-neighbor averaging and are
read at two sinks. With activity confined to two sources, the graph reduces to
$T=R_SW^2E_{ab}$, whose columns record source access at the readout. For the
noiseless subspace XOR task used in the archived experiments, both columns
being nonzero are necessary and sufficient for perfect Bayes accuracy. This
criterion is deliberately task-specific. It does not constrain arbitrary
GNNs, noisy decoded states, or finite-optimizer behavior.

The audit separates three levels that structural reports usually combine.
The first is graph structure: degrees, module counts, named motifs, protected
edges. The second is population information: whether a task label remains
decodable from the transferred sources. The third is finite trained behavior:
how a particular readout handles the surviving mixture. A common-relay count
is sufficient, but not necessary, for the population criterion on this task.
Finite accuracy can also vary widely among interventions that all preserve
ideal label information.

Our contributions are:
\begin{itemize}
  \item An operator-level audit for connectome interventions. It reports
        source visibility and a task-specific Bayes criterion beside the
        stated graph constraints.
  \item A matched MaleCNS counterexample. Two zero-common-relay, degree- and
        module-matched reconstructions have different source access; the
        high-versus-chance split is consistent with their archived fits.
  \item Converging tests of that distinction.
        All nonzero $k=8$ doses retain ideal information while finite accuracy
        rises, nulls can regain a lost task, executed readouts follow transfer
        geometry on visual labels, and a reverse-support image control stays
        negative.
\end{itemize}
